\section{Related Work}

\subsection{Flaky-test detection and classification}

A large line of work detects or predicts flakiness from execution
history~\cite{pinto2020understanding,alshammari2021flakeflagger,
lam2020idflakies}, vocabulary and static features~\cite{pinto2020understanding,
haben2023flakiness}, or dynamic behavioral features~\cite{fatima2022flakify};
more recently, large language models classify flaky
tests~\cite{kim2024large,yang2024predicting,fatima2023flakycat}. These tools
answer \emph{whether} a test is flaky, but do not reproduce the failure or
explain its mechanism. FlakeAgent is complementary: it takes a test known (or
suspected) to be flaky and produces the failing execution plus a grounded
diagnosis, derived from reproduced failures rather than static signals.

\subsection{Flaky-test reproduction}

iDFlakies~\cite{lam2019idflakies} pioneered automated reproduction of
order-dependent failures by executing whole test suites in randomized orders
and minimizing polluter--victim pairs; its companion plugin infrastructure
(testrunner~\cite{oei2020testrunner}) is the execution substrate we reuse for
our same-JVM double-run strategy. NonDex~\cite{gyori2016nondex} reproduces
implementation-dependent failures by perturbing hash and reflection
nondeterminism; our T12 strategy invokes it directly. FlakeRake~\cite{
bell2017flakerake} and Shaker~\cite{silva2017practical} inject sleeps and
perturb resources to expose timing sensitivity — the family of ideas behind
our sleep-injection and stress strategies. TDRepro~\cite{chen2024tdrepro}
specifically reproduces time-dependent failures by first learning timing
features from a \emph{given failure log} and then searching parameterized
stress levels to match it; it requires the failing log as input and covers
only TD. NIODebugger~\cite{shen2024nidebugger} localizes and fixes
non-idempotent outcomes but assumes the failure has already been observed.
FlakeRepro~\cite{alghmadi2023flakerepro} targets .NET and combines reruns
with fault localization, without an LLM component. What distinguishes
FlakeAgent is category agnosticism combined with cold start: no failure log,
no stack trace, no category hint — the agent reasons over the test source
and its own observations only, and covers all four categories with one loop.
We deliberately compare against a log-free baseline family (Section~\ref{sec:setup})
so that the comparison is apples-to-apples; TDRepro-style log-based tools are
inapplicable to the cold-start setting by construction.

\subsection{LLM agents for software engineering}

The agent paradigm — LLM in a loop with tools and observations — has produced
strong results on repository-level coding (SWE-bench~\cite{jimenez2024swebench}),
autonomous bug fixing~\cite{zhang2024agentless,zhang2024autocoderover}, and
debugging~\cite{zhang2024llamadebug,kang2023largelang}. These settings assume
the failure is reproducible on demand; flaky failures are different, because
the agent's central problem is \emph{making} the bug happen via
hypothesis-driven environment manipulation rather than code editing. Closest
to our work, TDRepro~\cite{chen2024tdrepro} couples an LLM pipeline to the
failing log to steer its stress search, and LDB~\cite{zhang2024llamadebug}
debugs by verifying runtime execution step by step — both assume the failure
has already been observed. We are, to our knowledge, the first to demonstrate
a closed-loop agent that reproduces failures from cold start and diagnoses
them with evidence evaluated against developer fixes, at scale, across
flakiness categories and across 2026-era frontier models. FlakeAgent also
follows the budgeted-agent design of~\cite{zhang2024agentless}: a hard budget
of six LLM calls and eight executions per test keeps the amortized cost in
the cents range (Section~\ref{sec:cost}), a precondition for adoption in
everyday CI debugging.
